\documentclass[10pt,a4paper]{amsart}
\usepackage[margin=19mm]{geometry}
\usepackage{amsmath,amssymb,mathtools}
\usepackage[scr=rsfs,cal=euler]{mathalfa}
\usepackage{xfrac}
\usepackage[unicode,hidelinks,bookmarks=false]{hyperref}
\newcommand{\ie}{i.e.\ }
\newcommand{\cf}{cf.\ }
\newcommand{\resp}{respectively\ }
\newcommand{\Subsection}{\subsection}
\providecommand{\nobreakdash}{\nobreak\hbox{-}}
\setlength{\emergencystretch}{2em}
\multlinegap0pt
\usepackage{amssymb}
\newtheorem{theorem}{Theorem}
\title{New binary optimal LCD codes using heuristic embedding}
\author{Haeun Lim, Junmin An, Jon-Lark Kim}
\date{Source: arXiv:2609.02096v1 (CC BY 4.0)}
\begin{document}
\maketitle

\noindent\emph{Source and licence.} New binary optimal LCD codes using heuristic embedding. Version arXiv:2609.02096v1 (CC BY 4.0), distributed by arXiv under the Creative Commons Attribution 4.0 International licence, http://creativecommons.org/licenses/by/4.0/, as recorded in the arXiv licence metadata for this exact version and shown on the abstract page https://arxiv.org/abs/2609.02096v1. Full licence notice: \texttt{licenses/arxiv-2609.02096-CC-BY-4.0.txt}.\par\smallskip

\noindent\textit{Editorial scope.} Counted unit: the article's theorem CMTQ-2019 (source0.tex lines 135--140). The article's complete original proof of it is delivered with it (source0.tex lines 164--183, one \texttt{proof} environment of 20 lines, which the article places away from the statement). No mathematics is written, repaired or re-derived: the statement and the proof are the article's own lines, in the article's own notation, and no retained line is substituted or re-flowed.  Retained uncounted as the source-local context the proof needs: lines 142--146.  Nothing was omitted from any retained span, so the package carries no omission record.  The retained lines cite 2 works, whose entries are transcribed verbatim from the article's own bibliography source in the e-print and delivered with short editorial labels recorded in the item's reference mapping.  No retained line contains a figure environment, an image-inclusion command or drawing code, so no image is delivered and the package declares no asset dependency.  Editorial formatting only, in addition: an AMS \texttt{amsart} preamble that restates the article's own declarations and macros used by the retained lines, namely the environments \texttt{theorem} on separate counters, together with the standard packages those lines need.  The retained environments print in this document as Theorem 1 in order of appearance, whereas the article prints the same items with the numbering of its own full document.  The complete native source of the article as distributed in the arXiv e-print for this exact version is bundled unchanged as \texttt{sources/209173/source0.tex}.
\begin{theorem}[\cite{LSL-2024}]\label{LCD-finite}
    If $2 \leq k \leq n$, then \[
    d_{LCD}(n,k) \leq \max \{d_{LCD}(n-1, k-1), d_{LCD}(n-2, k-2)\}.
    \]
\end{theorem}

\begin{theorem}[\cite{CMTQ-2019}]\label{CMTQ-2019}
    If $2 \leq k \leq n$, then
    \[
    d_{LCD}(n,k) \leq d_{LCD}(n, k-1).
    \]
\end{theorem}

\begin{proof}
The assertion follows by induction on $k$.  The initial values are exact,
and $d_{\mathrm{LCD}}(n,k)\le d_{g}(n,k)$ because every LCD
code is a linear code. Theorem~\ref{CMTQ-2019} gives the second candidate
in the minimum as
\[
d_{LCD}(n, k)\le d_{LCD}(n, k-1)\le U(n, k-1)
\]
for $6\le k\le 10$ and Theorem~\ref{LCD-finite} gives the third as
\[
\begin{aligned}
d_{LCD}(n,k)
&\leq \max\{d_{LCD}(n-1,k-1),
             d_{LCD}(n-2,k-2)\}\\
&\leq \max\{U(n-1,k-1),U(n-2,k-2)\}.
\end{aligned}
\]
Thus each of the three entries in the defining minimum is an upper bound
on $d_{LCD}(n,k)$, so their minimum is also an upper bound.
\end{proof}

\begin{thebibliography}{99}
\bibitem[CMTQ-2]{CMTQ-2019} \newblock C. Carlet, S. Mesnager, C. Tang and Y. Qi, \newblock New characterization and parametrization of LCD codes, \newblock IEEE Trans. Inform. Theory, 65(1) (2019), 39--49.
\bibitem[LSL-20]{LSL-2024} \newblock S. Li, M. Shi and H. Liu, \newblock Several constructions of optimal LCD codes over small finite fields, \newblock Cryptogr. Commun., 16 (2024), 779--800.
\end{thebibliography}
\end{document}
